\documentclass[11pt,a4paper]{article}
\usepackage[margin=1in]{geometry}
\usepackage{hyperref}
\usepackage{enumitem}
\usepackage{graphicx}
\usepackage{xcolor}
\usepackage{amsmath}

% -----------------------------------------------------
% bvraghav-tiet-ucs503p-article.sty
% -----------------------------------------------------

% Variable: Lab Instructor
\makeatletter \newcommand{\BvrSetLabInstructor}[1]{%
  \gdef\Bvr@LabInstructor{#1}}
\newcommand{\Bvr@LabInstructor}{Dr. Raghav %
  B. Venkataramaiyer}
\newcommand{\BvrLabInstructorValue}{\Bvr@LabInstructor}
\makeatother

% Add subtitle
\usepackage{titling}
\pretitle{\begin{center}\bfseries\fontsize{18bp}{18bp}\selectfont}
\makeatletter
\newcommand{\BvrSubtitle}[1]{%
  \posttitle{%
    \par\end{center}
    \begin{center}\large#1\end{center}
    \vskip 0.5em}%
}%
\makeatother

% Hints in section titles
\newcommand{\BvrHint}[1]{%
  {\normalfont\normalsize\itshape\hfill #1}}

% -----------------------------------------------------

\title{VeriHire -- Internship \& Placement Opportunity Credibility Platform}

\BvrSetLabInstructor{Dr. Jeelani}

\BvrSubtitle{%
  Project Proposal for UCS503P  \\
  Submitted to: \BvrLabInstructorValue \\ \bfseries
  Thapar Institute of Engineering and Technology%
}

\author{
  Geet Rawal, CSED%
  \thanks{Roll No: \texttt{1024160012}, %
    Email: \texttt{grawal\_be24@thapar.edu}}
  \and
  Nidhish, CSED%
  \thanks{Roll No: \texttt{1024160013}, %
    Email: \texttt{nnidhish\_be24@thapar.edu}}
  \and
  Parvesh Lamba, CSED%
  \thanks{Roll No: \texttt{1024160017}, %
    Email: \texttt{plamba\_be24@thapar.edu}}
}

\date{\today}

\begin{document}
\maketitle

\section{Higher-order goal%
  \BvrHint{}}

This project aims to improve the safety and reliability of internship and
placement discovery for students by reducing the uncertainty involved in
evaluating online opportunities. Students frequently encounter internship
and job opportunities through social media, messaging groups, professional
networking platforms, emails, and independent websites, where credibility
can be difficult to verify.

The immediate engineering objective is to develop a measurable and deployable
machine-learning-based software solution that can identify patterns associated
with credible and potentially suspicious internship and placement opportunities.
The system will provide students with an ML-based assessment that can assist
them in deciding whether additional verification is required before applying.

\section{Time-to-value%
\BvrHint{}}

\begin{itemize}[leftmargin=0.7cm]
    \item We will first implement the core workflow of opportunity submission,
    preprocessing, ML inference, and result visualization.

    \item We will develop the system through short iterations, allowing both
    the software workflow and machine learning pipeline to be tested continuously.

    \item Classical ML models will be implemented as baselines, followed by
    the proposed semantic-embedding and structured-feature MLP architecture.

    \item Model performance will be measured using precision, recall, F1-score,
    confusion matrices, and prediction latency.

    \item Success is defined by achieving a functional end-to-end ML workflow
    in the first iteration and improving the model and application through
    subsequent testing and feedback.
\end{itemize}

\section{Problem Statement}

Students increasingly discover internship and placement opportunities through
multiple online channels. However, they often lack a reliable and systematic
way to evaluate the credibility of these opportunities before applying.

Common pain points include:
\begin{itemize}[leftmargin=0.7cm]

    \item \textbf{Information fragmentation:} opportunity details are often
    distributed across job posts, company websites, messages, emails, and
    recruiter profiles.

    \item \textbf{Suspicious requirements:} some opportunities may request
    registration fees, deposits, paid training, or other forms of payment
    before selection.

    \item \textbf{Misleading claims:} unrealistic salary promises, guaranteed
    placements, urgency-based language, or exaggerated benefits can make
    suspicious opportunities difficult to identify.

    \item \textbf{Unstructured information:} internship descriptions contain
    textual, numerical, and categorical information that is difficult to
    evaluate manually.

    \item \textbf{Limited reporting:} students generally lack a centralized
    platform through which suspicious opportunities can be reported and
    information can be shared with others.

    \item \textbf{Manual verification:} students often have to independently
    investigate companies, recruiters, and job descriptions before deciding
    whether to apply.
\end{itemize}

\textbf{Why this matters now (contextual relevance):}
Students frequently make internship and placement decisions under time
constraints and with incomplete information. A machine-learning-based system
can learn patterns from labelled historical data and provide an initial
credibility assessment, reducing the effort required for manual evaluation.

\section{Proposed Solution%
\BvrHint{}}

\subsection{Overview}

Build a \textbf{web-based} ``Opportunity-to-Credibility'' platform with the
following components:

\begin{itemize}[leftmargin=0.7cm]

    \item \textbf{Opportunity submission portal:} students submit internship
    or placement information such as company, role, description, salary,
    location, recruiter information, and application details.

    \item \textbf{Data preprocessing:} submitted information is cleaned,
    normalized, and transformed into textual, numerical, and categorical
    features suitable for machine learning.

    \item \textbf{ML prediction engine:} the proposed hybrid model combines
    semantic text embeddings with structured opportunity features and uses a
    task-specific Multi-Layer Perceptron (MLP) to predict the opportunity's
    credibility/risk category.

    \item \textbf{Prediction interface:} the system displays the prediction,
    confidence/probability information, and supporting indicators to the
    student without presenting the result as absolute proof of fraud.

    \item \textbf{Community reporting:} students can report suspicious
    opportunities and provide additional information or evidence.

    \item \textbf{Admin/moderator dashboard:} authorized users can review
    opportunities, reports, verification decisions, and model predictions.

    \item \textbf{Analytics:} administrators can view statistics related to
    opportunities, predictions, reports, model performance, and system usage.
\end{itemize}

\subsection{Core workflow}

\begin{enumerate}[leftmargin=0.7cm]

    \item A student submits an internship or placement opportunity and
    receives an opportunity ID.

    \item The system validates and normalizes the submitted information.

    \item The opportunity text is prepared for semantic embedding while
    numerical and categorical attributes are encoded as structured features.

    \item A pretrained \textbf{all-MiniLM-L6-v2} Sentence Transformer generates
    a semantic representation of the opportunity text.

    \item The semantic embedding and encoded structured features are fused and
    passed to the task-specific MLP classifier.

    \item The model generates a predicted credibility/risk category and
    associated probability/confidence information.

    \item Supporting indicators are generated from relevant input features
    and the prediction is stored in the system.

    \item The result is displayed to the student. Students can report
    suspicious opportunities, while moderators can review reports and
    verification status.
\end{enumerate}

\subsection{Operational constraints for an educational setting}

\begin{itemize}[leftmargin=0.7cm]

    \item Keep the solution usable on common student devices through
    responsive web design.

    \item Use a suitable labelled dataset for supervised machine learning
    and maintain a clear separation between training, validation, and testing
    data.

    \item Use a pretrained sentence encoder as a feature extractor rather
    than training a large language model from scratch.

    \item Keep the task-specific MLP and inference pipeline lightweight enough
    to provide predictions within an acceptable response time.

    \item Avoid claiming absolute fraud detection; the system will provide
    an ML-based assessment to support student decision-making.

    \item Design the system for deployment using common web hosting,
    managed databases, and a lightweight Python ML inference service.
\end{itemize}

\section{Solution Approach%
\BvrHint{}}

This is an engineering course project with a machine learning component.
The focus will be on developing a complete, deployable ML-enabled software
system rather than solely optimizing a machine learning algorithm.

The overall pipeline will be:

\[
\text{Data Collection}
\rightarrow
\text{Data Cleaning}
\rightarrow
\text{Feature Engineering}
\rightarrow
\text{Model Training}
\rightarrow
\text{Evaluation}
\rightarrow
\text{Deployment}
\]

\subsection{Machine learning pipeline%
\BvrHint{}}

The problem will be formulated as a \textbf{supervised classification problem}.
The dataset will contain internship/job opportunity examples with labels
representing their credibility/risk category.

Potential features include:
\begin{itemize}[leftmargin=0.7cm]

    \item Internship/job description
    \item Salary or stipend information
    \item Payment-related terms and requirements
    \item Experience requirements
    \item Company information
    \item Recruiter/contact information
    \item Application method
    \item Location and employment type
    \item Relevant textual patterns
\end{itemize}

The project will use a two-stage modeling strategy. First, lightweight
classical models will be used as \textbf{baselines}. TF-IDF features may be
combined with Logistic Regression, Support Vector Machine, Random Forest,
or Gradient Boosting to establish reference performance.

The proposed final architecture will use a pretrained
\textbf{Sentence Transformer, all-MiniLM-L6-v2}, as a semantic feature
extractor. Its text embedding will be combined with encoded numerical and
categorical opportunity features. The resulting fused feature vector will
be passed to a task-specific \textbf{Multi-Layer Perceptron (MLP)}.

The proposed MLP architecture is:

\[
\text{Fused Features}
\rightarrow
128
\rightarrow
64
\rightarrow
32
\rightarrow
\text{Output}
\]

The output layer will represent the selected credibility/risk classes.
The exact hidden-layer configuration may be tuned using validation
experiments while keeping the architecture feasible for a semester-long
project.

Model comparison will use accuracy, precision, recall, F1-score, confusion
matrices, and inference latency. The final architecture will be selected
based on predictive performance, generalization, interpretability of the
results, and deployment feasibility.

\subsection{NLP and feature engineering%
\BvrHint{}}

Since internship and placement opportunities contain substantial textual
information, NLP techniques will be used to convert descriptions into
machine-readable representations.

The proposed NLP pipeline is:

\[
\text{Opportunity Text}
\rightarrow
\text{Text Cleaning}
\rightarrow
\text{all-MiniLM-L6-v2}
\rightarrow
\text{Semantic Embedding}
\]

Structured information will be processed separately:

\[
\text{Numerical/Categorical Features}
\rightarrow
\text{Encoding/Scaling}
\rightarrow
\text{Structured Feature Vector}
\]

The two representations will then be combined:

\[
\text{Semantic Embedding}
+
\text{Structured Features}
\rightarrow
\text{Feature Fusion}
\rightarrow
\text{MLP}
\rightarrow
\text{Prediction}
\]

TF-IDF will be retained as a lightweight baseline for comparison rather than
as the primary proposed architecture. This provides an empirical comparison
between classical text classification and the proposed semantic feature
fusion approach.

\subsection{Supporting indicators and responsible interpretation%
\BvrHint{}}

VeriHire will not present an ML prediction as a definitive statement that an
opportunity is fraudulent or genuine. Instead, the system will communicate
the predicted credibility/risk category together with supporting indicators
derived from the submitted opportunity.

Examples of indicators may include payment requirements, unusual application
methods, unrealistic compensation patterns, suspicious textual signals, or
other features learned and identified during the project. These indicators
are intended to help the student understand why additional verification may
be appropriate.

\subsection{Web architecture%
\BvrHint{}}

A modular architecture will consist of:

\begin{itemize}[leftmargin=0.7cm]

    \item \textbf{Frontend:} responsive React + Vite interface for students,
    moderators, and administrators.

    \item \textbf{Backend API:} Node.js and Express.js services responsible
    for authentication, opportunity management, database operations,
    reporting, and application logic.

    \item \textbf{ML service:} Python-based FastAPI service responsible for
    preprocessing, Sentence Transformer embedding generation, structured
    feature encoding, model loading, and MLP inference.

    \item \textbf{Database:} PostgreSQL database storing users, opportunities,
    predictions, reports, verification decisions, feedback, and related
    information.

    \item \textbf{Model artifacts:} versioned sentence encoder configuration,
    trained MLP parameters, feature encoders/scalers, and model metadata.
\end{itemize}

The ML service will remain independent from the main application backend,
allowing the trained model to be updated without major changes to the
frontend or core application.

\subsection{CI/CD and fast delivery enabler}

CI/CD will be used to:

\begin{itemize}[leftmargin=0.7cm]

    \item Automatically build the frontend and backend and run tests on
    every change.

    \item Test ML preprocessing, feature encoding, and prediction functions.

    \item Produce repeatable deployment artifacts for staging.

    \item Enable safe and frequent releases of incremental features.
\end{itemize}

\section{Evaluation Criterion%
\BvrHint{}}

We define success using measurable metrics associated with both the
machine learning model and the core application workflow.

\subsection{Primary evaluation metric}

\textbf{ML Classification F1-score:} the primary evaluation metric will be
the F1-score obtained on a held-out test dataset.

\begin{itemize}[leftmargin=0.7cm]

    \item \textbf{Target:} achieve an initial F1-score of approximately
    $0.80$ or above, subject to dataset quality and class distribution.

    \item \textbf{Attribution:} calculated by comparing model predictions
    against the known labels in the held-out test dataset.
\end{itemize}

\subsection{Secondary evaluation metrics}

\begin{itemize}[leftmargin=0.7cm]

    \item \textbf{Precision:} proportion of predicted suspicious
    opportunities that are actually suspicious according to the test labels.

    \item \textbf{Recall:} proportion of suspicious opportunities correctly
    identified by the model.

    \item \textbf{Inference latency:} time required by the ML service to
    preprocess an opportunity and return its prediction.

    \item \textbf{Model comparison:} performance difference between the
    classical baseline models and the proposed embedding + structured
    features + MLP architecture.

    \item \textbf{User clarity:} student-reported understanding of the
    prediction using a simple 1--5 feedback question.

    \item \textbf{Reliability:} availability of login, submission, and
    prediction functionality during the pilot.
\end{itemize}

\subsection{Pilot validation plan%
\BvrHint{}}

\begin{itemize}[leftmargin=0.7cm]

    \item Prepare a labelled dataset containing credible and potentially
    suspicious internship/job opportunities.

    \item Divide the dataset into training, validation, and test sets and
    compare the classical baselines against the proposed MLP architecture.

    \item Integrate the selected model into the web application and conduct
    testing with approximately 10--20 student users.

    \item Collect feedback regarding usability, prediction clarity, and
    system performance and use the results for subsequent improvements.
\end{itemize}

\section{Scalability%
\BvrHint{}}

The proposition should scale in both theoretical and practical senses.

\begin{enumerate}[leftmargin=0.7cm]

    \item \textbf{Scalable (theoretically):} The system should support growth
    in the number of users and opportunities by:

    \begin{itemize}[leftmargin=0.7cm]
        \item decoupling frontend, backend, ML service, and database,
        \item enabling pagination and indexing for common queries,
        \item using stateless API design,
        \item separating model inference from the main application,
        \item maintaining versioned ML models.
    \end{itemize}

    \item \textbf{Operationally deployable (practically):} The system should
    run on common cloud infrastructure:

    \begin{itemize}[leftmargin=0.7cm]
        \item managed PostgreSQL database,
        \item platform-hosted frontend and backend,
        \item Python/FastAPI ML inference service,
        \item containerized or platform-hosted deployment,
        \item CI/CD pipeline for staging and production.
    \end{itemize}

\end{enumerate}

\section{Engine availability heuristic}

A mature set of tools and frameworks is available to implement VeriHire as
a machine-learning-enabled web project:

\begin{itemize}[leftmargin=0.7cm]

    \item \textbf{React + Vite:} responsive frontend development.

    \item \textbf{Node.js + Express.js:} backend REST APIs and application
    logic.

    \item \textbf{PostgreSQL:} structured relational database.

    \item \textbf{Python, Pandas and NumPy:} data processing and preparation.

    \item \textbf{Scikit-learn:} baseline model training, preprocessing,
    evaluation, and comparison.

    \item \textbf{Sentence Transformers:} pretrained semantic text
    embeddings using \texttt{all-MiniLM-L6-v2}.

    \item \textbf{PyTorch:} implementation and training of the task-specific
    MLP classifier.

    \item \textbf{FastAPI:} serving the trained ML pipeline through an API.

    \item \textbf{Git and GitHub:} version control and collaborative
    development.

    \item \textbf{Testing and CI tools:} automated unit/integration testing
    and deployment pipelines.
\end{itemize}

We will use these mature components to focus on data quality, feature
engineering, model evaluation, software architecture, correctness, and
measurement rather than developing new infrastructure.

\section{Project scope and deliverables%
\BvrHint{}}

\subsection{Initial deliverable%
\BvrHint{}}

\begin{itemize}[leftmargin=0.7cm]

    \item Student registration and login.

    \item Internship/job opportunity submission and opportunity ID generation.

    \item Opportunity browsing and basic search.

    \item Dataset preparation, cleaning, and preprocessing pipeline.

    \item TF-IDF/classical ML baseline for initial performance measurement.

    \item Sentence Transformer embedding pipeline using
    \texttt{all-MiniLM-L6-v2}.

    \item Structured feature encoding and feature-fusion pipeline.

    \item Task-specific MLP classifier and evaluation pipeline.

    \item Prediction API using Python/FastAPI.

    \item Web interface for displaying the ML prediction and supporting
    indicators.

    \item PostgreSQL database for users, opportunities, predictions, and
    reports.

    \item CI: build, linting, and backend/ML unit tests.

    \item CD to staging with a single command or pipeline job.
\end{itemize}

\subsection{Subsequent deliverables}

\begin{itemize}[leftmargin=0.7cm]

    \item Hyperparameter tuning and comparison of baseline and MLP models.

    \item Improved NLP and structured feature engineering.

    \item Prediction interpretation and supporting indicators.

    \item Community reporting and moderator review.

    \item Company verification and evidence management.

    \item Advanced search and filtering.

    \item Notification layer using in-app notifications, with email as an
    optional later enhancement.

    \item Admin dashboard for model, opportunity, and moderation analytics.

    \item Model version management and improved scale readiness.
\end{itemize}

\section{Risks and mitigations}

\begin{itemize}[leftmargin=0.7cm]

    \item \textbf{Dataset quality:} poor or biased training data can reduce
    model performance. \textbf{Mitigation:} perform data cleaning,
    exploratory analysis, duplicate removal, and appropriate dataset
    splitting.

    \item \textbf{Class imbalance:} one class may contain significantly more
    examples than another. \textbf{Mitigation:} analyze class distribution
    and use class weighting or suitable resampling techniques where required.

    \item \textbf{Overfitting:} the model may perform well on training data
    but poorly on unseen opportunities. \textbf{Mitigation:} use validation
    and test datasets, cross-validation where appropriate, dropout or other
    regularization techniques, and early stopping where applicable.

    \item \textbf{False predictions:} the model may incorrectly classify
    legitimate or suspicious opportunities. \textbf{Mitigation:} evaluate
    precision, recall, F1-score, and confusion matrices and clearly present
    predictions as decision support.

    \item \textbf{Model drift:} suspicious opportunity patterns may change
    over time. \textbf{Mitigation:} design the system so that new labelled
    data can be used for future model retraining.

    \item \textbf{Data privacy:} user submissions may contain personal or
    company information. \textbf{Mitigation:} minimize collected data and
    apply authentication and access controls.

    \item \textbf{Operational issues in pilot:} ML service or deployment
    failures may affect demonstrations. \textbf{Mitigation:} deploy the
    inference service early, test the prediction API independently, and
    maintain a stable model version for demonstration.
\end{itemize}

\section{Summary}

This project proposes \textbf{VeriHire}, a deployable machine-learning-based
web platform designed to help students evaluate the credibility of
internship and placement opportunities.

The final proposed ML workflow is:

\[
\text{Opportunity Text}
\rightarrow
\text{all-MiniLM-L6-v2 Embedding}
\]

\[
\text{Structured Features}
\rightarrow
\text{Encoding/Scaling}
\]

\[
\text{Semantic Embedding}
+
\text{Structured Features}
\rightarrow
\text{Feature Fusion}
\rightarrow
\text{MLP}
\rightarrow
\text{Credibility/Risk Prediction}
\]

The project will use classical TF-IDF-based models as baselines and compare
them with the proposed semantic-embedding and structured-feature MLP
architecture. The final model will be selected using precision, recall,
F1-score, confusion matrices, inference latency, and deployment feasibility.

It meets the course criteria by emphasizing rapid iterations, modular
software architecture, machine learning integration, CI/CD, testing,
scalability, and measurable evaluation criteria.

The proposed technology stack consists of \textbf{React, Vite, Node.js,
Express.js, PostgreSQL, Python, Pandas, NumPy, Scikit-learn, Sentence
Transformers, PyTorch, FastAPI, and Git/GitHub}. The project is feasible for
a team of 3--4 students within one semester, with advanced explainability,
analytics, moderation, notifications, and model improvements introduced
incrementally.

The ultimate objective is to help students move beyond simply asking
``Is this internship genuine?'' and instead receive a data-driven
assessment that helps them make a more informed decision before applying.

\end{document}
